% 02-lc-quantization — brief recap; full derivation lives in HBP and the thesis
\section{Light-cone quantization}
\label{sec:lcq}

We follow the conventions of Ref.~\cite{Hornbostel:1988fb} throughout.
Light-cone coordinates are $x^\pm = x^0 \pm x^1$, with $x^+$ playing the role
of time; the conjugate momenta are $p^\mp$, and the mass-shell condition
$p^+ p^- = m^2$ makes $p^-$ the light-cone energy of a quantum carrying
longitudinal momentum $p^+ > 0$.  Positivity of $p^+$ is the structural
advantage of the scheme: no combination of quanta with strictly positive
longitudinal momenta can add up to the vacuum's $p^+ = 0$, so the Fock vacuum
is an exact eigenstate of the full Hamiltonian and every basis state is built
from quanta that belong to the hadron, not to vacuum fluctuations.  The
familiar equal-time picture, in which even a two-particle state mixes with
arbitrarily complicated vacuum structure, is replaced by one in which the
partition of a single conserved quantity organizes the entire basis.

Discretization proceeds by placing the theory in a box of length $2L$ in
$x^-$ with antiperiodic boundary conditions for the fermions, so that
momenta come in half-odd-integer units $k^+_n = 2\pi n/L$,
$n = \tfrac12, \tfrac32, \ldots$.  The dimensionless resolution
\begin{equation}
  K \;=\; \frac{L}{2\pi}\,P^+
  \label{eq:Kdef}
\end{equation}
is a positive integer (in code units we track $2K$, which is odd or even as
the parton number is), and a basis state is nothing but a colour-singlet
partition of $K$ among quanta.  At $K = 3$, for example, the available
momentum fractions are the partitions $(3)$, $(2,1)$ and $(1,1,1)$ --- to be
contrasted with the unbounded sets of positive and negative momenta an
equal-time discretization admits at the same volume.  Because the smallest
available momentum fraction is $x = k/K \sim 1/K$, the resolution doubles as
an infrared regulator of the $x \to 0$ endpoint, and the continuum limit
$L \to \infty$ at fixed $P^+$ is the single limit $K \to \infty$.  How
observables approach that limit --- and when the $1/K$ endpoint cutoff is the
dominant error --- is the subject of Secs.~\ref{sec:budget} and
\ref{sec:chiral}.

Everything else --- the mode expansions, the elimination of the constrained
degrees of freedom, and the construction of $P^-$ --- is unchanged from
Ref.~\cite{Hornbostel:1988fb}, and we do not repeat it.
